A partir de las limitaciones identificadas, se propusieron las siguientes verificaciones para ampliar el alcance del estudio.
\begin{enumerate}
\item Se recomendó extender el registro temporal del balance, especialmente en vacío, para cuantificar transitorios y almacenamiento de energía durante varios períodos.
\item Se propuso refinar la geometría del conductor y de la fuente, conservando sus dimensiones físicas, para examinar el cambio de las potencias absolutas y ampliar las comparaciones de malla del balance.
\item Se recomendó validar la relación de pérdida radial con posiciones o corridas independientes de las usadas para ajustar sus parámetros. Para evaluar penetración a través de materiales se propuso introducir interfaces y espesores definidos, con flujos incidente, reflejado y transmitido.
\item Se propuso calibrar la excitación mediante una potencia o un campo físico de referencia antes de convertir las potencias normalizadas a vatios. La evaluación de eficiencia de radiación requirió además una superficie de control que permitiera relacionar el flujo con la potencia suministrada a la antena.
\item Se recomendó conservar los perfiles numéricos de todas las resoluciones, junto con los parámetros del ajuste, sus réplicas y los registros de recursos. Esto permitió plantear una reproducción completa de los intervalos y de los criterios de comparación.
\end{enumerate}
